%% ACM Workshop Paper Template
%% Clinical Domain Pre-training vs. General Instruction-Tuning
%% Compiled with: pdflatex paper.tex; bibtex paper; pdflatex paper.tex; pdflatex paper.tex

\documentclass[sigconf,anonymous=false,review=false]{acmart}

%% Remove ACM-specific copyright block for course submission
\setcopyright{none}
\acmConference[Course Paper]{Graduate Seminar: AI in Healthcare}{2026}{University}
\acmYear{2026}

\usepackage{booktabs}
\usepackage{microtype}
\usepackage{graphicx}
\usepackage{xspace}
\usepackage[utf8]{inputenc}

%% Convenience macros
\newcommand{\medgemma}{MedGemma-4B-IT\xspace}
\newcommand{\gemma}{Gemma-3-4B-IT\xspace}
\newcommand{\pval}[1]{$p{=}{#1}$}
\newcommand{\sig}{$p{<}0.001$}

% -----------------------------------------------------------------------
\begin{document}

\title{Vendor Specialization Claims at 4B:\\ A Controlled Comparison of MedGemma and Gemma-3\\ on MIMIC-III MICU Nursing Tasks}

\author{Anonymous}
\email{anonymous@university.edu}
\affiliation{\institution{University}}

\begin{abstract}
Large language models (LLMs) marketed as pre-trained on clinical corpora are
widely assumed to outperform general-purpose models on clinical natural language
processing (NLP) tasks. We test this vendor claim directly for a single
4-billion-parameter checkpoint pair, comparing
\medgemma (Google's clinically pre-trained model) against
\gemma, a general instruction-tuned model of identical size and architecture,
on three MICU nursing tasks drawn from MIMIC-III: structured nursing note
generation (Sub-task~1), abnormal vital sign identification (Sub-task~2), and
admission summary generation (Sub-task~3). Using one record per unique hospital
admission ($n = 637$ independent records) and 13 evaluation metrics, the
clinically marketed checkpoint provides no consistent advantage. \gemma matches
or outperforms \medgemma on the majority of metrics, including abnormal value
recall ($\Delta = +0.119$, $p < 0.001$), AARP document structure
($\Delta = +0.024$, $p < 0.001$), and named mention of out-of-range metrics
in free-form summaries ($\Delta = +0.079$, $p < 0.001$). We caution that the
two checkpoints differ in alignment recipe as well as pre-training corpus, so
these results falsify the vendor's specialization claim for this checkpoint
pair without identifying the pre-training corpus as the causal driver. We
report the result as null-to-negative and hypothesise that, at 4B parameters,
instruction-following quality dominates domain vocabulary as a determinant of
structured clinical generation performance.
\end{abstract}

\keywords{clinical NLP, large language models, nursing documentation, MIMIC-III,
domain pre-training, benchmarking}

\maketitle

% -----------------------------------------------------------------------
\section{Introduction}
\label{sec:intro}

The deployment of artificial intelligence in clinical settings has been
accompanied by a pervasive marketing claim: that models pre-trained on medical
literature, clinical notes, and biomedical ontologies will outperform
general-purpose models on clinical language tasks. This framing is intuitively
appealing, since a clinician trained on medical texts should know more
medicine than a generalist, and it has driven considerable research effort in
domain-specific pre-training pipelines for healthcare NLP
\cite{lee2020biobert,singhal2023large}.

The practical consequence is that ``clinical pre-training'' on a model card
risks functioning as a credibility signal that outweighs empirical evidence
from the institution's own patient population and workflow when health systems
evaluate LLMs for clinical decision support or automated documentation.

We argue that this signal deserves controlled scrutiny on the specific
checkpoints health systems are likely to consider. Nursing documentation in
the intensive care unit is a stress-test where domain knowledge should matter
most: notes are highly structured (vital signs, medications, active problems,
assessments), grounded in observable measurements, and subject to clinical
accuracy requirements. If a vendor's clinically marketed checkpoint cannot
provide a measurable advantage here, the vendor's specialization claim
requires revision for that checkpoint.

This paper provides exactly that test for one checkpoint pair. We compare
\medgemma \cite{medgemma2024} (Google's 4B-parameter model marketed as
pre-trained on clinical and biomedical corpora) against \gemma
\cite{gemmateam2025gemma3}, a general instruction-tuned model of identical
size and architecture, on three tasks derived from MIMIC-III
\cite{johnson2016mimic} MICU nursing records. We emphasise that the two
checkpoints differ not only in pre-training corpus but also in alignment
recipe, instruction-tuning data, and possibly base-model snapshot; our design
therefore falsifies the vendor's specialization claim for this pair without
identifying the pre-training corpus as the sole causal driver. To eliminate
within-admission correlation, we sample one record per unique hospital
admission ($n = 637$ independent admission records):

\begin{enumerate}
  \item \textbf{Sub-task~1}: Generate a structured MICU nursing note in
    Assessment/Action/Response/Plan (AARP) format from 12~hours of structured
    EHR context.
  \item \textbf{Sub-task~2}: Identify abnormal vital signs and laboratory
    values from recent observations.
  \item \textbf{Sub-task~3}: Generate a free-form prose admission summary
    comparable to a hospital course narrative.
\end{enumerate}

Our results are unambiguous: \medgemma does not consistently outperform
\gemma. After reclassifying near-floor metrics with trivial absolute
differences as draws (Section~\ref{sec:results}), the general model is the
better performer on the majority of significant metrics. We report this as
an explicit null-to-negative result for this checkpoint pair and argue it
has direct implications for how healthcare organisations interpret vendor
specialization claims when evaluating 4B-parameter LLMs for clinical NLP
applications.

% -----------------------------------------------------------------------
\section{Related Work}
\label{sec:related}

\paragraph{Clinical pre-training benchmarks.}
\citet{lee2020biobert} demonstrated that BERT models pre-trained on PubMed
abstracts and clinical notes (BioBERT) improved named entity recognition
and relation extraction on biomedical benchmark datasets. However, these gains
were measured on tasks where training and test distributions closely matched
the pre-training corpus. Subsequent work by \citet{alsentzer2019publicly}
found that ClinicalBERT, pre-trained exclusively on MIMIC-III clinical
notes, improved over BioBERT on NER and natural language inference tasks,
yet underperformed on de-identification benchmarks where the test-set PHI
format differed systematically from the MIMIC pre-training distribution.
This result illustrates a broader pattern: domain pre-training benefits
are sensitive to corpus--task distribution alignment, not domain membership
alone. Our study
extends this line of inquiry to the instruction-following era, where
pre-training data composition interacts with instruction-tuning alignment
in ways that are not yet well characterised.

\paragraph{LLM evaluation on medical question-answering.}
\citet{singhal2023large} evaluated PaLM and Med-PaLM on medical licensing
examination questions (MedQA, MedMCQA, PubMedQA) and found that
domain-specific fine-tuning substantially improved performance on
knowledge-recall benchmarks. A critical distinction, however, is that
multiple-choice medical QA measures knowledge retrieval, not clinical
language \emph{generation}. Generating a coherent, grounded, structured
nursing note from a patient's vital signs requires capabilities (structured
output compliance, numerical grounding, instruction following) that are
distinct from encyclopaedic medical knowledge. Our evaluation targets this
generation regime rather than the QA regime.

\paragraph{Automated nursing documentation.}
\citet{rajkomar2018scalable} showed that EHR-grounded deep learning models
could predict clinical events from structured data with high accuracy, but
did not address free-text generation. More recently, \citet{nori2023capabilities} demonstrated that GPT-4 achieves
expert-level performance on medical licensing examinations and clinical
vignette question-answering, often matching or exceeding domain-specialised
models, suggesting that strong general instruction-following can be
competitive with domain fine-tuning even on clinical tasks. Our study provides a direct controlled
comparison at the 4B parameter scale, where computational resource constraints
make the choice between a domain model and a general model practically
meaningful for resource-limited healthcare settings.

% -----------------------------------------------------------------------
\section{Methodology}
\label{sec:methods}

\subsection{Dataset}

We use MIMIC-III \cite{johnson2016mimic}, a de-identified EHR database of
58,976 hospital admissions at Beth Israel Deaconess Medical Center. We
restrict to the Medical Intensive Care Unit (MICU) and filter
\texttt{NOTEEVENTS} for nursing progress notes (\texttt{category = `Nursing'}).
After ETL processing (see below), the dataset comprises 103,442 records across
three splits (train 71,885 / validation 16,169 / test 15,388). All experiments
use the validation split, sampling one record per unique hospital admission
($n = 637$ independent records).

\subsection{ETL Pipeline}

Raw MIMIC-III CSVs are processed using Polars
(lazy evaluation, chunk-safe for large joins). ITEMID harmonisation uses
the MIMIC-Extract variable map \cite{wang2020mimic}, which maps each measurement
identifier to a canonical variable name; only entries with verified status
are retained. Physiological plausibility filtering applies
per-variable outlier bounds from MIMIC-Extract \cite{wang2020mimic};
observations outside these bounds are flagged as out-of-range.
For each nursing note, we construct a \emph{FHIR-lite} record containing:

\begin{itemize}
  \item \textbf{Observations}: all vital sign readings from the 12~hours
    preceding the note's \texttt{charttime}, with out-of-range flags derived
    from the MIMIC-Extract physiological bounds \cite{wang2020mimic},
    capped at 1,100 entries at prompt-construction time.
  \item \textbf{Labs}: all laboratory values from the same 12-hour window,
    capped at 150 entries at prompt-construction time.
  \item \textbf{Medications}: all CareVue medication administration events
    from the same window, capped at 250 entries at prompt-construction time.
  \item \textbf{Active problems}: ICD-9 diagnoses linked to the admission.
  \item \textbf{Hospital course}: The ``\textsc{Brief Hospital Course}'' or
    ``\textsc{Hospital Course}'' section extracted from the admission's
    discharge summary via regex, used as the reference for Sub-task~3.
\end{itemize}

Records are retained only if the density gate $|\text{obs}| + |\text{labs}| +
|\text{meds}| \geq 5$ is satisfied. Temporal data uses ISO~8601 relative
offsets (e.g., \texttt{T-10h}) to mask MIMIC date-shifting artefacts.
The per-section caps (1,100 observations, 150 labs, 250 medications) were
set to keep prompts within the model context window; 99.1\% of admissions
fall entirely within these limits, and both models received identically
constructed inputs.

\subsection{Models and Inference}

Both models are served via vLLM \cite{kwon2023efficient} on a single NVIDIA
RTX~3090 (24~GB VRAM), with \texttt{tensor\_parallel\_size=1}. Sampling
parameters are identical across all runs:

\begin{itemize}
  \item Temperature: 0.30; top-$p$: 0.90; repetition penalty: 1.15.
  \item Max tokens: 1,024 (Sub-tasks~1 and 3); 512 (Sub-task~2).
\end{itemize}

The repetition penalty of 1.15 was necessary for \medgemma to suppress
degenerate looping: in an initial diagnostic run on the first 500 records
(predating the final n=637 corpus) without this penalty, 95\% of
\medgemma outputs were hard-truncated at the token limit with no structured
content. \gemma exhibited this behaviour in only 14 of those 500 records.
We applied the same penalty to both models for comparability, but did not
sweep penalties per model; we therefore cannot distinguish \medgemma's
performance ceiling at its optimal configuration from its performance under
the only configuration we found that produced any output (see
Section~\ref{sec:limits}).

\paragraph{Sub-task~1 prompt.} The system prompt instructs the model to act as
an ICU nurse completing a shift note. The user turn provides all FHIR-lite
fields as structured text and requests one complete AARP block per active
problem, preceded by a single de-identified few-shot example drawn from the
validation split. The same exemplar is used across all 637 prompts and is
verified excluded from the n=637 evaluation set by \texttt{row\_id}.
The ground truth is the original nursing note from \texttt{NOTEEVENTS}.
Full prompt text is given in Appendix~\ref{app:prompts}.

\paragraph{Sub-task~2 prompt.} The model receives only the most recent reading
per vital sign metric and is asked to list every abnormal value with its
clinical significance. Ground truth is derived by applying standard ICU
clinical reference ranges (e.g., Heart Rate 60--100\,bpm, SpO$_2 \geq 95$\%)
that are fixed in the inference engine; these are clinically motivated normal
ranges, distinct from the statistical outlier bounds used in the ETL
physiological plausibility filter. The full reference range table and prompt
are given in Appendix~\ref{app:prompts}.

\paragraph{Sub-task~3 prompt.} The system prompt adopts a physician persona
and requests a free-form prose Hospital Course narrative. The user turn
provides all FHIR-lite fields, identical in structure to Sub-task~1.
The reference is the extracted hospital course section. Full prompt text
is given in Appendix~\ref{app:prompts}.

\subsection{Evaluation Metrics}

\paragraph{Sub-task~1.}
\begin{itemize}
  \item \textbf{ROUGE-L} (F-measure, stemmed) against the original nursing note.
  \item \textbf{Clinical Entity F1} (aggregate, drug surface-form, vital
    surface-form) via \texttt{en\_core\_sci\_sm} scispaCy NER.
  \item \textbf{AARP Section Score}: proportion of active problems receiving a
    complete Assessment/Action/Response/Plan block within a 2400-character
    sliding window.
\end{itemize}

\paragraph{Sub-task~2.}
Binary precision, recall, and F1 over abnormal vital sign identification,
where a value is ``identified'' if its label appears in the generated text.

\paragraph{Sub-task~3.}
\begin{itemize}
  \item \textbf{ROUGE-L} against a \emph{relevance-filtered} hospital course
    (sentences whose key terms appear in \texttt{source\_facts} only).
  \item \textbf{BERTScore F1} (\texttt{distilbert-base-uncased}) against
    the same filtered reference.
  \item \textbf{Problem Coverage}: fraction of ICD-9 problem descriptions
    appearing as verbatim lowercased substrings in the generated text.
  \item \textbf{Out-of-Range Term Mention Rate} (formerly ``Key Lab
    Coverage''): fraction of out-of-range vital/lab \emph{labels} (e.g.,
    ``heart rate'', ``lactate'') present as lowercased substrings anywhere
    in the generated text. The metric does not check that the corresponding
    numeric value is reproduced or that the model identifies the value as
    abnormal; it measures named mention only.
  \item \textbf{Medication Name Mention Rate} (formerly ``Medication
    Accuracy''): fraction of administered medications whose first
    whitespace token (e.g., ``morphine'' from ``Morphine 2 mg administered
    via IV'') appears as a lowercased substring in the generated text. The
    metric does not check dose, route, or that the drug is described as
    administered.
\end{itemize}

We use the term \emph{mention rate} rather than ``coverage'' or ``accuracy''
to reflect that these metrics reward vocabulary overlap, not clinical
grounding or numerical fidelity.

All continuous metric comparisons use the paired Wilcoxon signed-rank test
(two-sided) with rank-biserial correlation $r$ as a measure of direction
consistency across paired comparisons.\footnote{Rank-biserial $r$ reports
the proportion of paired comparisons in which one model dominates the other;
it is a measure of \emph{direction consistency}, not effect magnitude. Values
near 1 indicate one model wins on nearly every pair, even when the absolute
mean difference is small.} Significance threshold: $\alpha = 0.05$. We do
not apply a multiple-comparisons correction; results with marginal absolute
deltas (Section~\ref{sec:results}) should therefore be interpreted with
caution.

\subsection{Pipeline Overview}

Figure~\ref{fig:pipeline} illustrates the end-to-end workflow from raw
MIMIC-III CSVs through evaluation.

\begin{figure*}[tp]
\centering
\includegraphics[width=\textwidth]{figures/academic/fig0_pipeline.pdf}
\caption{End-to-end pipeline. MIMIC-III raw CSVs are processed via a
Polars-based ETL into FHIR-lite records. Phase~2 serves both models via vLLM
on a single GPU. Phase~3 computes task-specific metrics and performs paired
statistical comparison.}
\label{fig:pipeline}
\end{figure*}

% -----------------------------------------------------------------------
\section{Results}
\label{sec:results}

All three sub-tasks were evaluated over $n = 637$ independent admission records
(one record per unique \texttt{hadm\_id}) drawn from the MIMIC-III validation
split. Both models received identical inputs; the only experimental variable
is the pre-training corpus.

\subsection{Sub-task 1: Structured Nursing Note Generation}

Table~\ref{tab:notes} reports corpus-level means and Wilcoxon test results.
ROUGE-L scores are near-floor for both models (0.079 vs.\ 0.076); although
the difference reaches significance, the absolute delta of $+0.003$ is below
any threshold of practical meaning and we treat this metric as a draw.
Aggregate clinical entity F1 is similarly near-floor (0.059 vs.\ 0.054)
and we likewise treat it as a draw; the recogniser
(\texttt{en\_core\_sci\_sm}) has known weak recall on ICU drug surface forms
and abbreviations, and the absolute scores indicate the metric is dominated
by recogniser noise rather than genuine model differences. Both metrics
also share the limitation that the ground-truth nursing note contains
bedside-observed content (physical examination, subjective patient reports)
not present in the structured EHR context, so neither metric measures
clinical accuracy directly.

The clinically meaningful differences favour \gemma: AARP Section Score
($\Delta = +0.024$, \sig) and Entity F1 (vitals)
($\Delta = +0.016$, \pval{0.017}).

The AARP result is striking but must be interpreted with care: the vast
majority of \medgemma outputs scored zero (no complete AARP block detected),
compared to a non-trivial fraction of \gemma outputs achieving structure.
\medgemma consistently failed to follow the four-header structure even after
the repetition penalty was applied. We did not sweep the repetition penalty
per model (Section~\ref{sec:limits}); we therefore cannot rule out that
\medgemma has a different inference configuration under which it produces
AARP-structured output. The result establishes that \medgemma fails AARP
compliance \emph{at the only configuration we found that produced any
non-truncated output for it}, not that it is incapable of AARP compliance
in principle.

\begin{table}[t]
\caption{Sub-task 1: Nursing Note Generation ($n = 637$ pairs). MG = MedGemma-4B-IT; G3 = Gemma-3-4B-IT. $\Delta$ = MG minus G3. Bold = winning model per row; $*$ = $p < 0.05$ (uncorrected). $r$ is rank-biserial direction consistency, not effect magnitude. Rows marked ``draw'' have significant $p$ but absolute $\Delta$ below practical-meaning threshold or are dominated by recogniser noise.}
\label{tab:notes}
\small
\begin{tabular}{@{}lrrrrrl@{}}
\toprule
\textbf{Metric} & \textbf{MG} & \textbf{G3} & \textbf{$\Delta$} &
  \textbf{$p$} & \textbf{$r$} & \textbf{Verdict} \\
\midrule
ROUGE-L            & \textbf{0.079} & 0.076 & $+$0.003 & $<$0.001$^*$ & 0.638 & draw \\
Entity F1          & \textbf{0.059} & 0.054 & $+$0.005 & $<$0.001$^*$ & 0.610 & draw \\
Entity F1 (drugs)  & 0.022 & \textbf{0.025} & $-$0.004 & 0.430 & 0.994 & n.s. \\
Entity F1 (vitals) & 0.058 & \textbf{0.074} & $-$0.016 & 0.017$^*$ & 0.975 & G3 \\
AARP Score         & 0.000 & \textbf{0.025} & $-$0.024 & $<$0.001$^*$ & 0.999 & G3$^\dagger$ \\
\bottomrule
\end{tabular}
\\[2pt]
{\footnotesize $^\dagger$ Penalty-configuration confound; see text.}
\end{table}

\subsection{Sub-task 2: Abnormal Value Identification}

Table~\ref{tab:abnormal} reports precision, recall, and F1 for identifying
abnormal vital signs. All three metrics are significant. \gemma
substantially outperforms \medgemma on recall ($\Delta = +0.119$, \sig),
F1 ($\Delta = +0.056$, \sig), and precision ($\Delta = +0.030$, \sig).

The Sub-task 2 metric counts a value as ``identified'' if its label appears
in the generated text; it does not require the model to explicitly flag the
value as high or low. A more verbose model that names more vital signs in a
preamble can therefore inflate its recall without performing the abnormal-
flagging task. We cannot rule out that part of \gemma's recall advantage
reflects greater verbosity rather than superior abnormal-value
identification; a stricter re-scoring requiring explicit
high/low flagging is left for future work.

Figure~\ref{fig:heatmap} shows per-vital miss rates. \medgemma systematically
misses Heart Rate (32\% miss rate) and Respiratory Rate (28\%); \gemma
reduces these miss rates to 1\% and 3\% respectively. MAP remains difficult
for both models (42\% MedGemma, 50\% Gemma-3), possibly because mean arterial
pressure is a derived value that may not always appear as a direct
\texttt{ITEMID} in CareVue.

\begin{table}[t]
\caption{Sub-task 2: Abnormal Value Identification ($n = 637$ pairs).
Bold = winning model per row; $*$ = $p < 0.05$ (uncorrected).}
\label{tab:abnormal}
\small
\begin{tabular}{@{}lrrrrr@{}}
\toprule
\textbf{Metric} & \textbf{MG} & \textbf{G3} & \textbf{$\Delta$} &
  \textbf{$p$} & \textbf{$r$} \\
\midrule
Precision & 0.370 & \textbf{0.400} & $-$0.030 & $<$0.001$^*$ & 0.805 \\
Recall    & 0.692 & \textbf{0.811} & $-$0.119 & $<$0.001$^*$ & 0.972 \\
F1        & 0.453 & \textbf{0.509} & $-$0.056 & $<$0.001$^*$ & 0.828 \\
\bottomrule
\end{tabular}
\end{table}

\begin{figure*}[tp]
  \centering
  \includegraphics[width=\textwidth]{figures/academic/fig3_vital_heatmap.pdf}
  \caption{Per-vital miss rates for Sub-task~2. Darker cells indicate higher
  miss rates. \medgemma misses Heart Rate, MAP, and Respiratory Rate at
  notably higher rates than \gemma. MAP remains difficult for both models.}
  \label{fig:heatmap}
\end{figure*}

\subsection{Sub-task 3: Admission Summary Generation}

\begin{figure*}[tp]
  \centering
  \includegraphics[width=\textwidth]{figures/academic/fig6_summary_task.pdf}
  \caption{Sub-task~3 results. (A) ROUGE-L violin plots. (B) BERTScore F1
  violin plots. (C) Clinical mention-rate bar chart. \medgemma achieves
  marginally higher lexical/semantic similarity; \gemma substantially
  outperforms on out-of-range term mention rate. $*$ = $p < 0.05$.}
  \label{fig:summary}
\end{figure*}

Table~\ref{tab:summary} reports all five Sub-task~3 metrics;
Figure~\ref{fig:summary} shows per-record distributions. The pattern reverses
partially: \medgemma achieves higher ROUGE-L and BERTScore F1, consistent
with its reported pre-training on clinical narrative text
\cite{medgemma2024}. However, \gemma names a higher fraction of out-of-range
metrics and administered medications: Out-of-Range Term Mention Rate
($\Delta = +0.079$, \sig) and Medication Name Mention Rate
($\Delta = +0.029$, \pval{0.003}). Problem Coverage is near-floor for both
models ($\approx$7\%) with no significant difference ($p = 0.672$),
reflecting the difficulty of verbatim ICD-9 description matching. We caution
that the mention-rate metrics reward vocabulary overlap, not numerical
fidelity or clinical correctness; the same verbosity-inflates-mentions
mechanism noted for Sub-task~2 likely contributes here.

\begin{table}[t]
\caption{Sub-task 3: Admission Summary Generation ($n = 637$ pairs). Filtered
reference = hospital course sentences with at least one term in
\texttt{source\_facts}. Bold = winning model per row; $*$ = $p < 0.05$
(uncorrected). Mention-rate metrics measure named-substring presence, not
numerical fidelity.}
\label{tab:summary}
\small
\begin{tabular}{@{}lrrrrr@{}}
\toprule
\textbf{Metric} & \textbf{MG} & \textbf{G3} & \textbf{$\Delta$} &
  \textbf{$p$} & \textbf{$r$} \\
\midrule
ROUGE-L                       & \textbf{0.110} & 0.101 & $+$0.009 & $<$0.001$^*$ & 0.743 \\
BERTScore F1                  & \textbf{0.739} & 0.732 & $+$0.007 & $<$0.001$^*$ & 0.808 \\
Problem Coverage              & 0.068 & \textbf{0.069} & $-$0.001 & 0.672 & 0.935 \\
Out-of-Range Mention Rate     & 0.235 & \textbf{0.314} & $-$0.079 & $<$0.001$^*$ & 0.964 \\
Medication Mention Rate       & 0.452 & \textbf{0.481} & $-$0.029 & 0.003$^*$ & 0.800 \\
\bottomrule
\end{tabular}
\end{table}

\subsection{Cross-task Summary}

Of 13 evaluated metrics across three sub-tasks, we classify two Sub-task~1
results (ROUGE-L, aggregate Entity F1) as draws despite reaching nominal
significance, on the grounds that the absolute deltas (\(+0.003\) and
\(+0.005\)) are below practical-meaning thresholds and the metrics are
near-floor or recogniser-limited. \medgemma retains a significant advantage
on 2 metrics (ROUGE-L and BERTScore F1 in Sub-task~3, both stylistic-
similarity measures with modest absolute gains). \gemma retains a significant
advantage on 7 metrics (Entity F1 vitals and AARP Score in Sub-task~1;
Precision, Recall, and F1 in Sub-task~2; Out-of-Range Term Mention Rate and
Medication Name Mention Rate in Sub-task~3). The remaining 4 metrics are
draws or not significantly different (Sub-task~1 ROUGE-L and aggregate
Entity F1 reclassified as draws; Sub-task~1 Drug Entity F1 and Sub-task~3
Problem Coverage not significant). The AARP and mention-rate results carry
the caveats discussed in the preceding sections.

% -----------------------------------------------------------------------
\section{Discussion}
\label{sec:discussion}

\subsection{Null-to-Negative Result Interpretation}

This study is an explicit null-to-negative result, and we consider that its
most important contribution. The common inference pattern, ``this model was
pre-trained on clinical data, therefore it should perform better on clinical
tasks,'' is not supported by our data. At the 4B parameter scale, with
identical architecture and inference parameters, the general instruction-tuned
model matches or exceeds the clinically pre-trained model on the majority of
clinically meaningful metrics.

We hypothesise two explanations. First, the tasks we evaluate (structured
output compliance, AARP format, precise numerical grounding, and instruction
following) are primarily \emph{capability} requirements, not \emph{knowledge}
requirements. A model that generates a correctly structured AARP block with
accurate vital sign values succeeds through instruction adherence and pattern
recognition, not through medical knowledge retrieval. The clinical pre-training
corpus likely contained extensive unstructured clinical narratives, which may
have conditioned \medgemma away from rigid template adherence.

Second, the \medgemma repetition pathology (looping under long structured
prompts) suggests that the clinical pre-training or alignment process may
have introduced a sensitivity to long, repetitive structured contexts that
the general model does not exhibit. This is a practical limitation: a model
that requires a repetition penalty to produce any coherent output is
operationally less reliable than one that does not.

\subsection{Sub-task~3 Mixed Pattern}

Sub-task~3 is the only task where \medgemma shows a significant advantage,
and the advantage is on \emph{stylistic} metrics (ROUGE-L, BERTScore), not on
named mention of input facts. This is consistent with the explanation above:
\medgemma's pre-training corpus reportedly includes clinical notes and
discharge summaries \cite{medgemma2024}, which share the prose style of the
hospital course reference text. The model has learned to generate
discharge-summary-like prose, but that prose names fewer of the specific
out-of-range metrics and medications from the structured input than \gemma's
output does. Because our mention-rate metrics measure substring presence
rather than numerical fidelity, we frame this as a difference in
\emph{vocabulary alignment with the input} rather than as clinical grounding;
nevertheless, in a clinical safety context, naming the relevant out-of-range
metric is a necessary (if not sufficient) condition for grounded
documentation.

\subsection{Limitations}
\label{sec:limits}

\paragraph{Independence by admission.}
This evaluation samples one record per unique hospital admission ($n = 637$
unique admissions). This design satisfies the paired independence assumption
of the Wilcoxon test and eliminates the within-admission ICD-9 correlation
present in note-level sampling. Results represent 637 statistically independent
patients.

\paragraph{ROUGE-L and Entity F1 floor effect.}
Sub-task~1 ROUGE-L is 0.076--0.079 for both models, near-floor; aggregate
Entity F1 is similarly low at 0.054--0.059.
The ground-truth nursing notes are written from direct bedside observation
(physical examination, subjective patient reports) that is entirely absent
from the structured EHR context. No model can recover this information.
Neither metric is a valid quality signal for this task; we retain them only
for comparability with prior work and reclassify them as draws.

\paragraph{Inference configuration confound.}
We applied an identical repetition penalty of 1.15 to both models because
preliminary diagnostics showed \medgemma degenerated into truncated loops
without it. We did not subsequently sweep the penalty per model. Each
model's performance ceiling at its own optimal inference configuration is
therefore unknown, and the AARP Score result in Sub-task~1 in particular
should be read as ``\medgemma fails AARP at the only configuration we found
that produced any non-truncated output for it,'' not as a configuration-
independent capability claim.

\paragraph{Sub-task 2 verbosity confound.}
The Sub-task 2 metric counts a vital sign as identified if its label appears
in the generated text, without requiring the model to flag the value as
abnormal. A more verbose model can therefore inflate recall by naming more
metrics in preamble. We did not re-score with a stricter explicit-flagging
rule.

\paragraph{Sub-task 3 mention-rate metrics.}
Out-of-Range Term Mention Rate and Medication Name Mention Rate measure
substring presence of metric labels and medication first-tokens
respectively. They do not check that numeric values are accurate, that
medications are described as administered, or that out-of-range values are
explicitly identified as abnormal. They reward vocabulary alignment with
the input, not clinical grounding in a strong sense.

\paragraph{BERTScore encoder.}
BERTScore F1 was computed with \texttt{distilbert-base-uncased}, a
general-domain distilled encoder that may over-reward stylistic fluency
relative to clinical specificity. A clinical encoder
(e.g., \texttt{Bio\_ClinicalBERT}) may yield different relative scores
between the two models on Sub-task~3.

\paragraph{Single ICU type and institution.}
All records are from the MICU at a single Boston academic medical centre.
Nursing note conventions, vital sign reference ranges, and medication
administration practices differ across ICU types (cardiac, surgical,
neonatal) and institutions. Generalisability is unknown.

\paragraph{Parameter scale.}
Both models are 4B parameters; this is simultaneously a design strength and a
limitation. The controlled comparison (identical size, identical architecture,
identical inference parameters) isolates pre-training corpus as the single
experimental variable, but the results speak only to this scale. Two distinct
scale effects remain unexplored: (1)~whether the \emph{absolute capability
floor} at 4B suppresses task performance for both models equally, masking
pre-training effects that emerge only at higher capacity; and (2)~whether the
\emph{relative ordering} between clinical and general pre-training is stable
across scales. We observed that \medgemma's primary failure mode (AARP
instruction non-compliance) is a general instruction-following deficit; this
class of failure tends to diminish at larger parameter counts for all models.
However, \gemma's larger variants (12B, 27B) also benefit from the same
scaling, so the relative advantage is not guaranteed to reverse. Evaluating
\medgemma-27B against \gemma-27B would be the natural follow-up.

\paragraph{Single checkpoint pair.}
We compare exactly two checkpoints. They differ in pre-training corpus, but
also in alignment recipe, instruction-tuning data, RLHF, and possibly base
model snapshot. Our design therefore falsifies the vendor's specialization
claim for this checkpoint pair without isolating the pre-training corpus as
the causal driver. Generalisation to other clinical-vs.-general pairs
(e.g., Meditron-7B vs.\ Llama-2-7B; BioMistral vs.\ Mistral-7B) is left for
future work.

\paragraph{No multiple-comparisons correction.}
We report uncorrected $p$-values across 13 metrics. Several of the marginal
results in Sub-task~1 (ROUGE-L, aggregate Entity F1) have absolute deltas
below practical-meaning thresholds; we have reclassified these as draws in
the cross-task summary, but readers applying a Holm or BH-FDR correction
should not be surprised if some borderline significances do not survive
adjustment.

% -----------------------------------------------------------------------
\section{Conclusion}
\label{sec:conclusion}

We evaluated \medgemma against \gemma on three MICU nursing tasks from
MIMIC-III, using one record per unique hospital admission ($n = 637$
independent records) and 13 metrics. For this checkpoint pair at the 4B
parameter scale, the clinically marketed model does not provide a consistent
performance advantage. After reclassifying near-floor Sub-task~1 metrics as
draws, \gemma significantly outperforms \medgemma on 7 of 13 metrics,
including AARP document compliance (subject to an inference-configuration
caveat), abnormal value recall and precision, and named mention of
out-of-range metrics and medications in summaries. \medgemma retains
significant but modest advantages on stylistic similarity in free-form prose
(ROUGE-L and BERTScore F1 in Sub-task~3). Because our design varies both
pre-training corpus and alignment recipe between the two checkpoints, we
cannot attribute this pattern to pre-training corpus alone.

These findings have a practical implication: for resource-constrained
healthcare settings evaluating this specific 4B-parameter checkpoint pair
for MICU documentation support, the evidence does not justify a preference
for \medgemma on the basis of vendor specialization claims.
We hypothesise, without proving, that the more important variable is
instruction-following quality, which is shaped by alignment training rather
than pre-training corpus composition.

\paragraph{Future work.}
Two extensions would strengthen these findings: (1) evaluate
\medgemma-27B against \gemma-27B to test whether the relative ordering is
stable across parameter scales;
(2) conduct a prospective human evaluation in which ICU nurses rate clinical
accuracy and usability of generated notes, which would provide the ground
truth that automated metrics cannot. We also identify the automated detection
of ICD-9 problem coverage as an open metric design problem: verbatim
description matching is too strict; clinical concept matching via UMLS
normalisation would be more informative.

% -----------------------------------------------------------------------
\clearpage
\bibliographystyle{ACM-Reference-Format}
\bibliography{references}

% -----------------------------------------------------------------------
\appendix

\section{Prompt Templates}
\label{app:prompts}

All prompts follow the chat format: a fixed system message sets the model
persona and constraints; a user message provides the structured patient
context assembled from the FHIR-lite record (Section~\ref{sec:methods}).
Placeholders in angle brackets represent fields populated at runtime.

\subsection*{Sub-task 1: Nursing Note Generation}

\textbf{System:}
\begin{quote}
\small
You are an experienced MICU nurse writing an end-of-shift progress note.
Use clinical shorthand. Do not invent medications, procedures, or diagnoses
not present in the context. Do not include a document header.
\end{quote}

\textbf{User:}
\begin{quote}
\small
\textit{[Admission context: patient identifier, care unit, admitting
diagnosis, admit time offset]}\\
\textit{[Active problems: ranked ICD-9 diagnosis list]}\\
\textit{[Observations: vital signs from the 12-hour window, with
out-of-range flags, capped at 1,100 entries]}\\
\textit{[Labs: laboratory values from the 12-hour window, capped at 150 entries]}\\
\textit{[Medications: CareVue administration events, capped at 250 entries]}\\
\textit{[Procedures: CareVue procedure events]}\\

\medskip
\textbf{Example note (for format reference only; do not copy content)}

\smallskip
Impaired Skin Integrity\\
Assessment:\\
All skin surfaces as documented in flowsheets. Some sloughing on buttocks,
pt stooling loose bilious green BM.\\
Action:\\
Buttocks cleansed, antifungal cream and adaptic applied. Flexiseal ordered
and inserted.\\
Response:\\
Buttock areas free of fecal contamination overnoc, areas appear slightly
less irritated.\\
Plan:\\
Cleanse and treat impaired skin areas, protect open areas from friction and
fecal contamination.

\medskip
Renal failure, acute (ARF)\\
Assessment:\\
CRRT as ordered, no difficulties with circuit. UO via foley 10--15 cc clear
yellow.\\
Action:\\
Removed 1--200 cc/h. Ca and K gtts titrated per labs.\\
Response:\\
Attained fluid removal goal, neg 2700 cc yesterday. Hemodynamically stable.\\
Plan:\\
Fluid removal goal 2.5--3 L/day. Monitor lytes q6 and titrate repletion
gtts as ordered.

\medskip
\textbf{Task}\\
Write the nursing progress note. One AARP block per active problem.
Synthesise the data; do not restate it verbatim.
\end{quote}

The few-shot example is drawn from a de-identified MIMIC-III nursing note
in the validation split and is included in the user turn so the model
treats it as a concrete format demonstration.

\subsection*{Sub-task 2: Abnormal Value Identification}

\textbf{System:}
\begin{quote}
\small
You are an experienced MICU nurse reviewing a patient's vital signs.
Identify which values are clinically abnormal for an ICU adult patient.
For each abnormal value: (1) state whether it is high or low,
(2) give the normal adult range with units,
(3) name the clinical concern in one sentence.
If a value is within normal limits, do not mention it.
Be concise. Use clinical shorthand.
\end{quote}

\textbf{User:}
\begin{quote}
\small
\textbf{Patient Context}\\
Admitting diagnosis: \textit{[admit diagnosis]}\\
Care unit: \textit{[care unit]}

\medskip
\textbf{Most Recent Vital Signs}\\
\textit{[One line per metric: name, value, unit, time offset]}

\medskip
\textbf{Task}\\
Identify and explain any clinically abnormal vital signs above.
\end{quote}

Only the nine metrics with hardcoded ICU reference ranges are included
(Heart Rate 60--100\,bpm; Systolic BP 90--140\,mmHg; Diastolic BP
50--90\,mmHg; MAP 65--105\,mmHg; Respiratory Rate 12--20\,breaths/min;
SpO$_2 \geq 95$\%; Temperature 97.7--99.5\,\textdegree F;
CVP 2--12\,mmHg; Glucose 70--180\,mg/dL). These same ranges define the
ground-truth labels used for evaluation and are independent of the
statistical outlier bounds applied during ETL.

\subsection*{Sub-task 3: Admission Summary Generation}

\textbf{System:}
\begin{quote}
\small
You are an experienced MICU attending physician writing the Hospital Course
section of a discharge summary. Synthesise the structured ICU data provided
into a concise problem-oriented narrative. For each active problem: describe
the clinical findings, key investigations, interventions, and patient
response. Use clinical language. Be specific: cite abnormal values,
medications, and procedures by name. Do not use bullet points or AARP
headings; write flowing prose.
\end{quote}

\textbf{User:}
\begin{quote}
\small
\textit{[Admission context, active problems, observations, labs,
medications, procedures; identical fields to Sub-task 1]}

\medskip
\textbf{Task}\\
Write the Hospital Course narrative for this 12-hour ICU window.
One paragraph per active problem.
Synthesise the data; do not restate it verbatim.
\end{quote}

\end{document}
